%! AP Statistics — Unit 1, Lesson 1.2 — Lesson Plan
\documentclass[10pt]{article}
\usepackage[margin=0.6in]{geometry}
\usepackage{apstats-boxes}
\usepackage{enumitem}
\usepackage{multicol}
\usepackage{amsmath,amssymb}
\usepackage{array}
\usepackage{tabularx}

\newcolumntype{C}[1]{>{\centering\arraybackslash}p{#1}}
\newcolumntype{Y}{>{\centering\arraybackslash}X}
\setlist[itemize]{leftmargin=1.2em,itemsep=0.35em,topsep=0.35em}
\setlist[enumerate]{leftmargin=1.4em,itemsep=0.35em,topsep=0.35em}
\setlength{\parindent}{0pt}
\setlength{\parskip}{3pt}

\begin{document}

\begin{center}
    {\Large\bfseries AP Statistics: 2026--2027\ (55 minutes) \\[4pt]
    {\small\bfseries Unit 1: Exploring One-Variable Data \qquad
     Lesson 1.2: The Language of Variation --- Variables}}
\end{center}

\begin{tcolorbox}[colback=sky,colframe=navy,boxrule=0.9pt,arc=2mm,
    left=2mm,right=2mm,top=1.5mm,bottom=1.5mm]
    \textbf{Primary Objective:} Students will classify variables as \textit{categorical} or
    \textit{quantitative}, and further distinguish quantitative variables as \textit{discrete} or
    \textit{continuous}. Students will articulate why variable type determines which statistical
    methods are appropriate (Big Idea: VAR).
\end{tcolorbox}

\vspace{4pt}

\begin{skillbox}[Priority Ideas \& Skills]{goldbox}
\begin{minipage}[t]{0.48\linewidth}
    \textbf{AP Skill 2 --- Data Analysis}
    \begin{itemize}
        \item Correctly identify whether a variable is categorical or quantitative in context.
        \item Recognize that categorical data produce counts/proportions; quantitative data
              produce numerical summaries.
        \item Identify variable type before choosing a display or summary statistic.
    \end{itemize}
\end{minipage}
\hfill
\begin{minipage}[t]{0.48\linewidth}
    \textbf{Key Understandings}
    \begin{itemize}
        \item A number is not automatically quantitative (ZIP code, jersey number).
        \item The \textit{average test}: ``Does averaging these values give a meaningful
              result?'' separates categorical from quantitative.
        \item Discrete vs.\ continuous matters for probability models in Unit 4.
    \end{itemize}
\end{minipage}
\end{skillbox}

\begin{skillbox}[{Vocabulary, Concepts, \& Theorems}]{greenbox}
\begin{minipage}[t]{0.48\linewidth}
    \begin{itemize}
        \item \textbf{Variable} --- a characteristic that can take different values
              for different individuals
        \item \textbf{Categorical variable} --- places an individual into one of
              several groups or categories; no meaningful arithmetic
        \item \textbf{Quantitative variable} --- takes numerical values for which
              arithmetic makes sense
    \end{itemize}
\end{minipage}
\hfill
\begin{minipage}[t]{0.48\linewidth}
    \begin{itemize}
        \item \textbf{Discrete quantitative variable} --- countable values with
              gaps (e.g., number of siblings)
        \item \textbf{Continuous quantitative variable} --- can take any value in
              an interval; measured, not counted (e.g., height, time)
        \item \textbf{Distribution of a variable} --- what values a variable takes
              and how often; developed fully in Lessons 1.3--1.10
    \end{itemize}
\end{minipage}
\end{skillbox}

\begin{skillbox}[Activate Prior Knowledge \& Spiral Review]{sky}
\begin{minipage}[t]{0.48\linewidth}
    \textbf{Warm-Up Worksheet (5 min)}\\
    Students examine a small dataset from Lesson 1.1 (commute times study) and
    answer: ``Which columns record words or labels? Which record numbers?
    Can you average the mode of transport? Can you average the commute time?''
    Use responses to motivate today's classification.\\[6pt]
    \textbf{Topics Connected}
    \begin{itemize}
        \item[\rule{0.7cm}{0.4pt}] Variable types introduced in Lesson 1.1
        \item[\rule{0.7cm}{0.4pt}] What ``individual'' and ``variable'' mean
        \item[\rule{0.7cm}{0.4pt}] Preview: choosing appropriate displays (L1.3--1.9)
    \end{itemize}
\end{minipage}
\hfill
\begin{minipage}[t]{0.48\linewidth}
    \textbf{Connections to Current Learning}
    \begin{itemize}
        \item Lessons 1.3--1.4 display \textit{categorical} variables with bar
              charts and two-way tables.
        \item Lessons 1.5--1.9 display and summarize \textit{quantitative}
              variables with dotplots, histograms, and boxplots.
        \item Knowing variable type now prevents method-mismatch errors
              all year.
    \end{itemize}
\end{minipage}
\end{skillbox}

\begin{skillbox}[Hook \& Lesson]{sky}
\begin{minipage}[t]{0.48\linewidth}
    \textbf{Hook (4--5 min)}\\
    Project a data table with 5 variables about students: name, grade level, GPA,
    favorite color, hours of sleep. Ask students to sort the variables into two
    groups \textit{without} being told the group names. Debrief: What was the
    sorting rule? Use their language to introduce \textit{categorical}
    vs.\ \textit{quantitative}.
\end{minipage}
\hfill
\begin{minipage}[t]{0.48\linewidth}
    \textbf{Lesson Flow (15 min)}
    \begin{enumerate}
        \item Define categorical and quantitative; demonstrate the ``average test.''
        \item Work through 8 examples as a class, including traps:
              ZIP code, jersey number, 1--5 Likert rating.
        \item Introduce discrete vs.\ continuous with concrete anchors:
              \textit{number of pets} (count, gaps) vs.\
              \textit{weight of a pet} (measured, any value).
        \item Students propose one variable of each type from their own lives;
              share out briefly.
    \end{enumerate}
\end{minipage}
\end{skillbox}

\begin{skillbox}[Explicit Instruction \& Active Monitoring]{sky}
\begin{minipage}[t]{0.48\linewidth}
    \textbf{Direct Instruction (10 min)}
    \begin{itemize}
        \item Stress trap variables: area code, student ID, Likert scale.
              Key question: ``Does averaging jersey numbers give a meaningful result?''
        \item Show that the same \textit{question} can yield different variable
              types: ``How do students rate the cafeteria?'' as a 1--5 scale
              \textit{or} as ``satisfied / unsatisfied.''
        \item Decision flowchart on board: Is it a word/label?
              $\to$ Categorical. Is it a count?
              $\to$ Discrete. Otherwise $\to$ Continuous.
    \end{itemize}
\end{minipage}
\hfill
\begin{minipage}[t]{0.48\linewidth}
    \textbf{Active Monitoring --- Mini-Whiteboard Round}
    \begin{itemize}
        \item Display a variable; students hold up \textbf{C} (categorical),
              \textbf{QD} (quantitative discrete), or \textbf{QC}
              (quantitative continuous).
        \item Watch for: ordinal ratings marked QD without justification.
        \item Prompt: ``Can you find the average? Does that average
              \textit{mean} something in context?''
        \item Use activity worksheet Part 1 (Problems 1--6) as the
              whiteboard-round prompts before the group card sort.
    \end{itemize}
\end{minipage}
\end{skillbox}

\begin{skillbox}[Group Work \& Differentiation]{redbox}
\begin{minipage}[t]{0.48\linewidth}
    \textbf{Card Sort Activity (8--10 min)}\\
    Groups receive the activity worksheet with 12 variable cards and sort
    them into three columns: \textbf{Categorical} /
    \textbf{Quantitative Discrete} / \textbf{Quantitative Continuous}.
    After sorting, groups compare with neighbors and resolve disagreements.
    Debrief 2--3 contested cards whole-class (jersey number, star rating,
    Likert scale).
\end{minipage}
\hfill
\begin{minipage}[t]{0.48\linewidth}
    \textbf{Differentiation}
    \begin{itemize}
        \item \textit{Support}: Reference the decision flowchart on the
              activity sheet while sorting.
        \item \textit{Extension}: For each categorical variable on the card sort,
              propose a re-design that would make it quantitative instead
              (and vice versa).
        \item \textit{Discussion}: Debrief ``number of stars in a review'' ---
              is it categorical or quantitative? What does context tell us?
    \end{itemize}
\end{minipage}
\end{skillbox}

\begin{skillbox}[Individual Work \& Assessment]{redbox}
\begin{minipage}[t]{0.48\linewidth}
    \textbf{Exit Ticket (5 min)}\\
    Five music-themed variables --- students classify each and justify:
    \begin{enumerate}
        \item Favorite music genre
        \item Number of songs on a playlist
        \item Duration of the longest song (in minutes)
        \item Country of birth of the artist
        \item Temperature at the outdoor concert venue at noon
    \end{enumerate}
    \textcolor{slate}{\small Completeness graded; justification required.}
\end{minipage}
\hfill
\begin{minipage}[t]{0.48\linewidth}
    \textbf{AP-Style Check}\\
    \textit{``A survey records annual income (in dollars) and job type
    (teacher, engineer, other) for each respondent. Which statement is
    correct?''}
    \begin{enumerate}[label=(\Alph*),itemsep=2pt]
        \item Annual income is categorical.
        \item Job type is quantitative.
        \item Annual income is quantitative; job type is categorical.
        \item Both variables are categorical.
    \end{enumerate}
    \textcolor{slate}{\small Answer: C. Grade on classification + justification.}
\end{minipage}
\end{skillbox}

\begin{skillbox}[Reinforcement \& Extension]{goldbox}
\begin{minipage}[t]{0.48\linewidth}
    \textbf{Homework / Practice}
    \begin{itemize}
        \item Given an NBA player dataset (name, team, points per game, position,
              games played, jersey number): classify all 6 variables with
              justification.
        \item Three AP-style multiple-choice problems on variable identification,
              including at least one trap variable.
    \end{itemize}
\end{minipage}
\hfill
\begin{minipage}[t]{0.48\linewidth}
    \textbf{Extension / Preview}
    \begin{itemize}
        \item \textit{Extension}: Look up a dataset on the CDC or Census website;
              identify 3 categorical and 3 quantitative variables and note which
              are discrete vs.\ continuous.
        \item \textit{Preview}: Lesson 1.3 --- once we know a variable is
              categorical, how do we summarize and display it? Think about
              frequency and relative frequency.
    \end{itemize}
\end{minipage}
\end{skillbox}

\end{document}
